\documentclass[sigconf,screen,review]{acmart}
\setcopyright{none}
\copyrightyear{2027}
\acmYear{2027}
\acmConference[WWW '27]{The Web Conference 2027}{May 10--14, 2027}{Dublin, Ireland}
\acmDOI{}
\acmISBN{}

\title{RuleShift-Web: A Replayable Workbench for Auditing Agent Memory under Changing Web Policies}
\author{Chih-Kai Wang}
\affiliation{\institution{National Taipei University of Education}\city{Taipei}\country{Taiwan}}

\begin{document}
\begin{abstract}
Web agents can retain a policy response that is still fresh under its HTTP
cache lifetime even after a site begins enforcing a revised rule. A successful
revalidation can also become obsolete before a consequential transaction
commits. RuleShift-Web is a local research workbench that makes these two
boundaries observable and replayable. It publishes versioned retail policies
as real HTTP resources, records source-linked agent evidence, enforces an
application-level policy revision check at transaction commit, and exposes
the sequence of responses, publications, attempted actions, and state changes
in an interactive trace. A frozen study evaluates five evidence strategies
on 3,200 executions over paired tasks from two open models, plus 1,600
decisions from a deterministic controller. In this one-relevant-source
harness, source-dependent validation matches all-source conditional
validation while avoiding four policy reads per task; unconditional refresh
has higher observed accuracy for both models. The workbench
therefore exposes the cost and limits of memory validation rather than
claiming a universally superior controller. Visitors can reproduce policy
changes and inspect saved development-model traces without a conference
network connection.
\end{abstract}

\ccsdesc[500]{Information systems~Web applications}
\ccsdesc[300]{Computing methodologies~Intelligent agents}
\keywords{Web agents, agent memory, HTTP caching, dynamic policies, reproducible evaluation}
\maketitle

\section{Motivation and scope}
An agent may read a return policy at time $t_1$, carry a derived action to
$t_2$, and encounter a policy publication between those times. Three objects
then diverge: the response held by the client, the representation selected by
the HTTP origin, and the policy enforced at transaction commit. A cache-fresh
response is not necessarily current in the site's business process; an HTTP
\texttt{304 Not Modified} validates a selected representation at the time of
that request, not an action performed later \cite{rfc9110}. These distinctions
are difficult to inspect when a benchmark collapses all policy access into an
abstract tool call.

Prior work already studies obsolete memory and downstream behavior
\cite{stale}, memory evolution in changing environments \cite{evoarena},
dependency-scoped validation before external actions \cite{planfence},
and agent grounding in environmental evidence \cite{envtrust}.
RuleShift-Web does not claim to invent stale-memory detection,
dependency-scoped validation, conditional HTTP requests, or transaction
revision checks. Its candidate contribution is an experimental instrument in
which visitors can schedule policy publication, acquire or reuse evidence,
and submit a transaction while inspecting a fixed commit-time guard and its
event trace. PlanFence already checks declared action dependencies against their
owners \cite{planfence}; here the visible objects are Web responses and validators, cache age,
and a separate application revision check. Their separation is the
workbench's subject, not a new validation protocol. Table~\ref{tab:prior}
positions that narrower probe against the cited studies.

\begin{table}[t]
\caption{Prior evaluation emphasis versus the interactive Web probe.}
\label{tab:prior}
\small
\begin{tabular}{lp{0.59\columnwidth}}
\toprule
Work & Main inspectable object \\
\midrule
STALE \cite{stale} & Implicit conflict and revised belief \\
EvoArena \cite{evoarena} & Memory evolution across changing environments \\
EnvTrust \cite{envtrust} & Environmental evidence and agent outcome \\
PlanFence \cite{planfence} & Plan lineage and action-time owner validation \\
This demo & HTTP representation, held cache evidence, and separate commit authority \\
\bottomrule
\end{tabular}
\end{table}

\section{System design}
\subsection{A policy site with separate authority surfaces}
The prototype runs a FastAPI origin on a local host. Policies for discounts,
shipping, returns, fulfillment, and address changes are served as HTTP
resources with ETags and a 120-second cache lifetime. An unchanged conditional GET
returns \texttt{304}; a relevant publication returns a new representation and
validator. A publication of an unrelated policy leaves the relevant validator
unchanged. A \texttt{200} alone does not prove that a representation changed.
Reverting A$\rightarrow$B$\rightarrow$A restores the representation's
original validator.
The application-level revision is also content-derived: after this reversion,
an old A revision can commit an action that is valid under current A. The
workbench tests current-rule consistency, not whether any intervening
publication occurred.

Agent-facing endpoints expose public policy, order inspection, and retail
operations. Researcher endpoints publish rules and expose traces; they
never enter agent prompts. Each transaction carries an expected policy
revision. The origin compares it with the active rule under the same lock
as the state mutation. A mismatch returns \texttt{409} without mutation.
This application-level guard is distinct from HTTP \texttt{If-Match} on
the policy resource. Eligibility is checked after the revision match, so
the trace separates stale attempts from current but prohibited actions by
their response error fields; both can have status \texttt{409}.

\subsection{Evidence and replay}
An evidence record retains its source URL, acquisition time, validator,
content revision, policy body, dependent source, and last successful
validation time. The controller can reuse a remembered response, refresh it
unconditionally, validate with \texttt{If-None-Match}, or check only sources
relevant to the pending action. A transport or HTTP error leaves verification
unresolved; it is never treated as a successful validation. If no HTTP
validator is available, a full GET and content comparison replace conditional
revalidation. The prompt builder
passes only the public task and evidence to a model. Research snapshots,
future publications, and scoring data remain outside model input.

The workbench records each policy read, publication, order inspection, and
transaction attempt with its response and resulting state hash. A verifier
reconstructs a run from a clean initial state and checks the resulting events
and state. The local web page lets a visitor publish policies, choose an
evidence check, submit an operation, and inspect the timeline. It uses an
application-managed evidence cache; explicit network checks bypass browser
cache storage so each intended request reaches the origin. A static bundle
contains a 41-event example and seven recorded purchase interventions, so
trace inspection and branch selection do not require conference Wi-Fi. The
offline branches display previously verified deterministic event records. A separately
labelled set of eight local Qwen development traces exposes raw model output,
evidence history, rejections, and final state without generating fresh model
decisions.

\section{Hands-on demonstration}
A visitor starts with policy P1 permitting a discount and reads its HTTP
representation. The visitor then publishes P2, which removes eligibility,
while leaving the client-held P1 response in place. Reusing that response
produces an attempted transaction with the old revision; the origin rejects
it without mutation. A conditional GET instead receives P2, after which the
visitor can run a non-LLM public-evidence check that denies without sending a
transaction. This interface does not generate a live model decision; saved
model decisions are inspected separately. An unconditional reload also uses
one GET; when unchanged it transfers a body rather than returning \texttt{304}.
Publishing an unrelated
shipping rule demonstrates that a discount validator remains unchanged.

In a second branch, the visitor validates a permissive policy, then publishes
a restrictive version just before commit. The transaction is again rejected
under the active revision. After a bounded refresh, the controller can reach
the current decision. Returning to P1 demonstrates that identical policy
content can regain the same validator. The interface shows the held evidence,
request and response, transaction result, committed state, and replay record.
The same origin provides return, shipment, and address-change examples. For a
conference venue, the intended deployment is a laptop-bound local service
with a static trace fallback and an interactive poster walkthrough. A
read-only table in both modes compares the frozen five-strategy results;
it is labelled separately from live user interventions and saved model
development traces. A prepared supplementary review bundle contains the offline
page, verified traces, local-origin code, manuscript, and audit summaries.

\section{Evaluation and limitations}
\subsection{Frozen protocol}
We generated four held-out change scenarios: no change, relevant policy
change, unrelated policy change, and A$\rightarrow$B$\rightarrow$A
reversion. Each has ten distinct eight-task streams; the same 320 task
positions are reused across strategies and models. The five compared
strategies are memory retrieval without proactive validation, unconditional
refresh of the relevant source, conditional validation of all five sources,
a prior selective-retry controller, and the proposed source-dependent
conditional validation. Each task declares one relevant policy source;
conditional-all checks five sources but passes the same relevant evidence to
the model. Initial site state, public observations, tasks, and action budget
are paired across strategies. An in-process FastAPI TestClient drives the
origin, with scheduled publications before each task's evidence check. The evaluated models are
Qwen3-4B-Instruct-2507 and Phi-4-mini-instruct at pinned revisions, using
greedy float16 CUDA inference on one NVIDIA L4. An exact Markdown-fence
adapter for Phi was selected on development data before the held-out protocol
was frozen; Qwen uses a strict parser. A deterministic public-evidence
controller runs the same streams as a strong non-LLM reference.

The primary outcome requires a correct final decision \emph{and} matching
transaction state. We preserve wrong actions, rejected transactions, raw
model output, response-contract failures, HTTP requests, token counts, source evidence,
and replayable origin events. The audit independently reconstructs the
world and checks every result against the frozen stream schedule. Paired
accuracy differences use 40 eight-task streams, stratified by scenario,
with 10,000 bootstrap draws and two-sided 95\% percentile intervals.
The 3,200 model executions and 1,600 separate control decisions completed
with no failed arms. The common server guard prevented invalid commits;
this does not imply that the agents avoided invalid attempts.

\begin{table}[t]
\caption{Correct final decision and state, out of 320 tasks per cell.}
\label{tab:results}
\small
\begin{tabular}{lrrr}
\toprule
Strategy & Qwen & Phi & Rule controller \\
\midrule
Retrieval & 273 & 141 & 300 \\
Prior selective & 274 & 141 & 300 \\
Always refresh & \textbf{287} & \textbf{223} & \textbf{320} \\
Conditional all & 284 & 195 & \textbf{320} \\
Source dependent & 284 & 195 & \textbf{320} \\
\bottomrule
\end{tabular}
\end{table}

\begin{table}[t]
\caption{Total agent HTTP requests over 320 tasks per cell, including actions.}
\label{tab:requests}
\small
\begin{tabular}{lrr}
\toprule
Strategy & Qwen & Phi \\
\midrule
Retrieval & 737 & 555 \\
Prior selective & 794 & 555 \\
Always refresh & 1,017 & 890 \\
Conditional all & 2,300 & 2,149 \\
Source dependent & 1,020 & 869 \\
\bottomrule
\end{tabular}
\end{table}

\subsection{What the traces show}
Source-dependent validation ties all-source conditional validation in
accuracy for both models. Conditional-all makes 1,800 policy reads per
model, versus 520 for source-dependent validation: the saving of four
reads per task follows from checking one declared source instead of five. Total
requests also include action attempts (Table~\ref{tab:requests}); they are
not a cost-adjusted success metric. Relative to retrieval,
source-dependent validation's paired accuracy gain is 3.44 percentage
points for Qwen (95\% interval 1.56--5.31) and 16.88 for Phi
(11.88--22.19). Always refresh has higher observed accuracy: 287 versus
284 for Qwen and 223 versus 195 for Phi. The source-minus-always intervals
are $[-1.88,0.00]$ and $[-16.88,-2.50]$ percentage points, respectively;
Qwen's interval includes zero. Source-dependent and always-refresh each
make 520 policy reads. Phi's 21 fewer total requests under source-dependent
validation correspond to 21 fewer attempted transactions, with lower
completion. Request count alone does not establish greater efficiency.

The no-LLM controller scores 320/320 with all three proactive validation
strategies and 300/320 with retrieval or prior selective retry. This
separates controller mistakes from avoidable memory staleness. Qwen's three
source-dependent errors relative to always refresh are address-age
judgments with identical policy body and revision but different evidence
acquisition timestamps. The model's prompts and decisions are retained;
these differences do not show an incorrect HTTP validator. Phi produces
1,626 raw parser rejections, many due to Markdown fences. After the
predeclared adapter, 685 \emph{response-level} contract failures remain,
including retry responses; all have contradictory \texttt{attempt} and
\texttt{reason} fields, not malformed JSON. In proactive Phi arms,
always-refresh has 97 final contract rejections versus 125 for
source-dependent validation; no accepted final decision is incorrect in
those arms. Thus the observed gap characterizes this configured
model--prompt--parser pipeline, not a demonstrated freshness advantage.
The adapter asymmetry also prevents a pure cross-model quality ranking.

The live origin was tested for unchanged \texttt{304}, changed
\texttt{200}, stale-commit \texttt{409}, irrelevant publication, reversion,
commit-time races, and source failures. Saved traces can be re-executed
against a clean initial state, checking HTTP responses and final state
hashes; the offline page only loads the previously verified records.
Seven live-origin sequences were previously reported as passing. In an
additional Safari run, the exported three-event trace (read \texttt{200},
publish \texttt{200}, purchase \texttt{409}) was independently replayed
with an unchanged final state. These live-socket and browser checks are
separate from the in-process formal matrix; that matrix does not test a
publication between validation and commit. The study uses binary decisions,
declared single-source dependencies, fixed publication schedules, a
synthetic retail site, and two small models. Its intervals describe
generated-stream variation rather than reliability on real commerce
websites. The local
conference deployment must still be checked on the presentation laptop.

\section{Ethical use of data}
The prototype uses synthetic policies and orders. It involves no human
participants, personal customer records, payment credentials, or real
transactions. Evaluated agents are confined to the controlled site routes;
researcher controls and scoring records are not agent observations. The
interactive service is intended to run locally rather than expose an
unauthenticated public endpoint.
Generative AI tools assisted with implementation, analysis, and manuscript
drafting. The sole human author is responsible for verifying the reported
evidence, citations, and claims before submission.

\begin{thebibliography}{9}
\bibitem{evoarena} J. Xu et al. 2026. EvoArena: Tracking Memory Evolution for Robust LLM Agents in Dynamic Environments. \emph{arXiv:2606.13681}. \url{https://arxiv.org/abs/2606.13681}.
\bibitem{stale} H. Chao, Y. Bai, R. Sheng, T. Li, and Y. Sun. 2026.
STALE: Can LLM Agents Know When Their Memories Are No Longer Valid?
\emph{arXiv:2605.06527}. \url{https://arxiv.org/abs/2605.06527}.
\bibitem{planfence} E. Chen, S. Wang, and C. G. Brinton. 2026.
Fresh Memory, Stale Plans: Dependency-Scoped Validation for Distributed
LLM-Agent Memory. \emph{arXiv:2609.03340}.
\url{https://arxiv.org/abs/2609.03340}.
\bibitem{envtrust} S. Sheng, Z. Wang, and L. Zhou. 2026.
When Agents Overtrust Environmental Evidence: An Extensible Agentic Framework
for Benchmarking Evidence-Grounding Defects in LLM Agents.
\emph{arXiv:2605.08828}. \url{https://arxiv.org/abs/2605.08828}.
\bibitem{rfc9110} R. Fielding, M. Nottingham, and J. Reschke. 2022.
HTTP Semantics. RFC 9110. \url{https://www.rfc-editor.org/rfc/rfc9110.html}.
\end{thebibliography}
\end{document}
